\documentclass[12pt]{article}
\usepackage[utf8]{inputenc}
\usepackage[T1]{fontenc}
\usepackage{hyperref}
\title{Vilela Language Resources for K-12 Curriculum Development}
\author{Compiled by Claude Sonnet 5 (claude-sonnet-5) for Andr\'as Kornai}
\date{}
\begin{document}
\maketitle

\section{Introduction}
Vilela (ISO 639-3 \texttt{vil}) is a language isolate of Chaco province, Argentina. Sources disagree on its status: some call it extinct, and a 2013 field note reports two semi-speakers (a brother and sister, aged 74 and 70 in 2008). Shift to Spanish was nearly complete by the early 1970s. It is sometimes grouped with Lule, but Zamponi (2008) treats the two as separate isolates.

\section{Primary Educational Resources}
No primers, textbooks or school programmes for Vilela were found.

\section{Linguistic and Reference Materials}
\begin{itemize}
\item Vilela language (Wikipedia): \url{https://en.wikipedia.org/wiki/Vilela_language}
\item Lule--Vilela languages (Wikipedia): \url{https://en.wikipedia.org/wiki/Lule%E2%80%93Vilela_languages}
\item Numerals page for Vilela, MPI-EVA Chaco numerals project: \url{https://lingweb.eva.mpg.de/channumerals/Vilela.htm}
\item Endangered Languages Project entry: \url{https://www.endangeredlanguages.com/fr/node/103280}
\end{itemize}
No dedicated Vilela dictionary or grammar was found in search results. The 1732 Machoni grammar mentioned in one listing covers Lule and Tonocot\'e, not Vilela.

\section{Religious and Community Resources}
None found.

\section{Digital and Audio Resources}
None found.

\section{Curriculum Development Context}
A revitalization curriculum would rest on scarce, mostly fieldwork-derived data. Specialist publications (for example by Zamponi or Viegas Barros) should be located and checked directly; they were not verified here.
\end{document}